\section*{Capítulo \ref{ch:g11:seq} --- Sequências: um primeiro curso}

\begin{solution}{exo:g11:seq:1}
$u_n = \frac{n}{n+1}$: $ u_1 = \frac12$, $ u_2 = \frac23$,
$ u_3 = \frac34$.

$ u_0 = 5$, $ u_{n+1} = 3u_n - 2$: $ u_1 = 13$, $ u_2 = 37$, $ u_3 = 109$.

$ u_n = (-1)^n n $: $ u_1 = -1$, $ u_2 = 2$, $ u_3 = -3$.
\end{solution}

\begin{solution}{exo:g11:seq:2}
$ u_{10} = 7 + 10 \times (-3) = -23$ e $ u_{25} = 7 - 75 = -68$.

Para $(v_n)$: $ v_8 = v_3 + 5d $ dá $26 = 11 + 5d $, então $ d = 3$; então
$ v_0 = v_3 - 3d = 11 - 9 = 2$.
\end{solution}

\begin{solution}{exo:g11:seq:3}
$ u_8 = 5 \times 2^8 = 1280$.

Para $(v_n)$: $ v_4 = v_2\, q^2$ dá $48 = 12 q^2$, então $ q^2 = 4$ e
$ q = 2$ (os termos são positivos). Então
$ v_0 = \frac{v_2}{q^2} = \frac{12}{4} = 3$.
\end{solution}

\begin{solution}{exo:g11:seq:4}
$1 + \dots + 500 = \frac{500 \times 501}{2} = 125\,250$.

$4 + 7 + \dots + 61$ é \omterm{def:g11:seq:arithmetic}{aritmética} com $ d = 3$ e
$\frac{61 - 4}{3} + 1 = 20$ termos: soma
$20 \times \frac{4 + 61}{2} = 650$.

$1 + \frac12 + \dots + \frac{1}{2^{10}}$ é \omterm{def:g11:seq:geometric}{geométrico} com $ q = \frac12$
e $11$ termos:
$\frac{1 - (1/2)^{11}}{1 - 1/2} = 2\left(1 - \frac{1}{2048}\right)
= \frac{2047}{1024}$.
\end{solution}

\begin{solution}{exo:g11:seq:5}
$ u_{n+1} - u_n = 4(n+1) - 1 - 4n + 1 = 4$: \omterm{def:g11:seq:arithmetic}{aritmética} com $ d = 4$.

$\frac{v_{n+1}}{v_n} = \frac{2^{n+1}}{3^{n+2}} \cdot \frac{3^{n+1}}{2^n}
= \frac23$: \omterm{def:g11:seq:geometric}{geométrico} com $ q = \frac23$.

$ w_0 = 0$, $ w_1 = 2$, $ w_2 = 6$: as diferenças $2$ e $4$ diferem, então
não \omterm{def:g11:seq:arithmetic}{aritmética}; $\frac{w_1}{w_0}$ nem sequer está definido, e os rácios
$\frac{w_2}{w_1} = 3 \neq \frac{w_3}{w_2} = 2$: nenhum.
\end{solution}

\begin{solution}{exo:g11:seq:6}
Os tamanhos das linhas são \omterm{def:g11:seq:arithmetic}{aritmética}: primeiro termo $16$, diferença $2$. O último
(20ª) linha tem $16 + 19 \times 2 = 54$ assentos. O total é
$20 \times \frac{16 + 54}{2} = 700$ assentos.
\end{solution}

\begin{solution}{exo:g11:seq:7}
Depois $ n $ horas a população está $500 \times 2^n $. Às 20h, $ n = 8$:
$500 \times 256 = 128\,000$ bactérias. Nós precisamos $500 \times 2^n >
10^6$, ou seja\ $2^n > 2000$: desde $2^{10} = 1024$ e $2^{11} = 2048$, o
população ultrapassa primeiro um milhão depois $11$ horário integral, às 23h.
\end{solution}

\begin{solution}{exo:g11:seq:8}
$ c_2 = 1.002 \times 100 + 100 = 200.20$ e
$ c_3 = 1.002 \times 200.20 + 100 \approx 300.60$. As diferenças
$ c_2 - c_1 = 100.20$ e $ c_3 - c_2 \approx 100.40$ não são iguais, então
$(c_n)$ não é \omterm{def:g11:seq:arithmetic}{aritmética}; as proporções $\frac{c_2}{c_1} = 2.002$ e
$\frac{c_3}{c_2} \approx 1.50$ também não são iguais, então não é
\omterm{def:g11:seq:geometric}{geométrico}. (Recursões mistas ``multiplicar e depois adicionar'' como esta são
resolvido pelo \omterm{def:g11:seq:sequence}{sequência} auxiliar truque de \cref{exo:g11:seq:11}.)
\end{solution}

\begin{solution}{exo:g11:seq:9}
$ u_{n+1} - u_n = (n+1)^2 - 8(n+1) - n^2 + 8n = 2n - 7$: negativo para
$ n \leq 3$, positivo para $ n \geq 4$. Então $(u_n)$ diminui para
$ u_4 = 16 - 32 = -16$, depois aumenta: não é monotônico.

$(v_n)$ tem termos positivos e
\[
\frac{v_{n+1}}{v_n} = \frac{3^{n+1}}{(n+1)!} \cdot \frac{n!}{3^n}
= \frac{3}{n+1},
\]
o que é $> 1$ para $ n \leq 1$, $= 1$ para $ n = 2$, e $< 1$ para
$ n \geq 3$: o \omterm{def:g11:seq:sequence}{sequência} aumenta até $ v_2 = v_3 = \frac92$, então
diminui.
\end{solution}

\begin{solution}{exo:g11:seq:10}
O primeiro $ n $ termos são $ u_0, \dots, u_{n-1}$, com $ u_0 = 3$ e
$ u_{n-1} = 3 + 4(n-1) = 4n - 1$. A soma deles é
\[
n \times \frac{3 + (4n-1)}{2} = n(2n + 1) = 903,
\]
então $2n^2 + n - 903 = 0$. Aqui $\Delta = 1 + 4 \times 2 \times 903 =
7225 = 85^2$, e $ n = \frac{-1 + 85}{4} = 21$ (o negativo \omterm{def:g11:quad:discriminant}{raiz} é
rejeitado). Verificar: $21 \times 43 = 903$.
\end{solution}

\begin{solution}{exo:g11:seq:11}
\emph{1.} $ u_1 = 5$, $ u_2 = 9$, $ u_3 = 17$: cada termo é um a mais que
$4, 8, 16$, sugerindo $ u_n = 2^{n+1} + 1$.

\emph{2.} Com $ v_n = u_n - 1$:
\[
v_{n+1} = u_{n+1} - 1 = 2u_n - 1 - 1 = 2(u_n - 1) = 2v_n,
\]
então $(v_n)$ é \omterm{def:g11:seq:geometric}{geométrico} com proporção $2$ e primeiro mandato
$ v_0 = u_0 - 1 = 2$.

\emph{3.} Por isso $ v_n = 2 \times 2^n = 2^{n+1}$ e
$ u_n = v_n + 1 = 2^{n+1} + 1$, confirmando a conjectura. (O número $1$
subtraído em $ v_n $ é o \omterm{pb:g10:functions:1}{ponto fixo} de $ x \mapsto 2x - 1$; o mesmo
ideia reaparece para $ u_{n+1} = au_n + b $ em \cref{ch:g12:seq}.)
\end{solution}

\begin{solution}{pb:g11:seq:1}
\textbf{1.} $ h_1 = 1$, $ h_2 = 3$, $ h_3 = 7$.

\textbf{2.} Para mover o disco maior, o $ n $ discos acima dele
deve primeiro migrar para a indexação sobressalente ($ h_n $ movimentos); o grande disco
cruzes ($1$ mover); o $ n $ os discos devem então subir de volta ao topo
disso ($ h_n $ movimentos): $ h_{n+1} = 2h_n + 1$.

\textbf{3.} $ v_{n+1} = h_{n+1} + 1 = 2h_n + 2 = 2v_n $:
\omterm{def:g11:seq:geometric}{geométrico} de proporção $2$ com $ v_1 = 2$, então $ v_n = 2^n $ e
$ h_n = 2^n - 1$.

\textbf{4.} $2^{64} - 1 \approx 1.8 \times 10^{19}$ segundos;
dividindo por $3 \times 10^7$ segundos por ano:
sobre $6 \times 10^{11}$ anos --- seiscentos bilhões de anos,
quarenta vezes a idade do universo. Os monges podem levar
pausas para café.

\textbf{5.} Em qualquer solução legal, considere o disco inferior
primeiro movimento: naquele instante o outro $ n $ os discos devem estar todos assentados
o único pino restante (pelo menos $ h_n $ se move para levá-los até lá),
e após o último movimento do disco inferior todos eles devem voltar
em cima dele (pelo menos $ h_n $ mais): qualquer solução precisa de pelo menos
$2h_n + 1$ se move. A recorrência é tanto um piso quanto um
teto: $2^n - 1$ é ideal.

\textbf{6.} $1, 1, 2, 3, 5, 8, 13, 21, 34, 55, 89, 144$.

\textbf{7.} Não \omterm{def:g11:seq:arithmetic}{aritmética} ($2 - 1 = 1$ mas $3 - 2 = 1$,
$5 - 3 = 2$: as diferenças mudam); não \omterm{def:g11:seq:geometric}{geométrico}
($\frac21 = 2$ mas $\frac32 = 1.5$). \omterm{def:g11:func:monotone}{Aumentando}: para
$ n \geq 2$, $ F_{n+1} - F_n = F_{n-1} > 0$.

\textbf{8.} $ F_k = F_{k+2} - F_{k+1}$, então
\[
\sum_{k=1}^{n} F_k = (F_3 - F_2) + (F_4 - F_3) + \dots +
(F_{n+2} - F_{n+1}) = F_{n+2} - F_2 = F_{n+2} - 1 .
\]
Para $ n = 6$: $1 + 1 + 2 + 3 + 5 + 8 = 20 = F_8 - 1 = 21 - 1$.

\textbf{9.} $ F_k F_{k+1} - F_{k-1} F_k = F_k (F_{k+1} -
F_{k-1}) = F_k \cdot F_k = F_k^2$; somando telescópios para
$ F_n F_{n+1} - F_1 F_0$ (com $ F_0 = 0$): a soma dos quadrados é
$ F_n F_{n+1}$. Para $ n = 4$: $1 + 1 + 4 + 9 = 15 = F_4 F_5 =
3 \times 5$.

\textbf{10.} $ F_5 F_3 - F_4^2 = 5 \times 2 - 9 = 1$;
$ F_6 F_4 - F_5^2 = 8 \times 3 - 25 = -1$;
$ F_7 F_5 - F_6^2 = 13 \times 5 - 64 = 1$: alternando
$\pm 1$. Este \omterm{def:g10:numbers:interval}{intervalo} entre $ F_{n+1} F_{n-1}$ e
$ F_n^2$ é exatamente a unidade quadrada ganha ou perdida do mágico:
cortando um $ F_n \times F_n $ quadrado em pedaços remontados como um
$ F_{n+1} \times F_{n-1}$ retângulo deve criar ou engolir um
unidade --- a lasca.

\textbf{11.} $ F_{n+2} = F_{n+1} + F_n \geq F_n + F_n = 2F_n $
(o \omterm{def:g11:seq:sequence}{sequência} aumenta): a cada dois índices, pelo menos um
duplicando --- crescimento pelo menos \omterm{def:g11:seq:geometric}{geométrico} de proporção $\sqrt2$ por
índice.

\textbf{12.} $1.5$; $1.667$; $1.6$; $1.625$; $1.615$;
$1.619$; $1.618$; $1.618$. Se $ r_n \to L $: de
$ F_{n+2} = F_{n+1} + F_n $, dividindo por $ F_{n+1}$:
$ r_{n+1} = 1 + \frac{1}{r_n}$, então $ L = 1 + \frac1L $, ou seja\
$ L^2 = L + 1$: $ L = \varphi = \frac{1 + \sqrt5}{2}$, o dourado
proporção de \cref{pb:g10:algebra:1}. Os coelhos se multiplicam em
ouro.

\textbf{13.} $ v_{n+1} = u_{n+1} - \ell = a u_n + b - \ell $;
desde $\ell = a\ell + b $, isso é $ a(u_n - \ell) = a v_n $:
\omterm{def:g11:seq:geometric}{geométrico} de proporção $ a $. Por isso $ v_n = a^n v_0$ e
$ u_n = a^n (u_0 - \ell) + \ell $.

\textbf{14.} \omterm{pb:g10:functions:1}{Ponto fixo}: $\ell = 1.01\ell - 300$ dá
$\ell = 30\,000$. Então
$ d_n = 1.01^n (10\,000 - 30\,000) + 30\,000
= 30\,000 - 20\,000 \times 1.01^n $.

\textbf{15.} $ d_n \leq 0$ requer $1.01^n \geq 1.5$:
$1.01^{40} \approx 1.489$, $1.01^{41} \approx 1.504$: o
$41$ O primeiro pagamento liquida a dívida (e é um pouco menor do que
$300$). Total reembolsado: pouco menos $41 \times 300 = 12\,300$
euros --- o $10\,000$ custo emprestado sobre $2\,300$ euros de
interesse.

\textbf{16.} \omterm{pb:g10:functions:1}{Ponto fixo} $\ell = \frac{1000}{1 - 1.02} =
-50\,000$, então $ p_n = 1.02^n \times 100\,000 - 50\,000$. Depois
$10$ anos: $1.02^{10} \approx 1.219$:
$ p_{10} \approx 71\,900$ habitantes.

\textbf{17.} $\frac{1000 \times 1001}{2} = 500\,500$; e
$2^{20} - 1 = 1\,048\,575$.

\textbf{18.} De $7$ para $502$ em etapas de $5$:
$\frac{502 - 7}{5} + 1 = 100$ termos; soma
$= 100 \times \frac{7 + 502}{2} = 25\,450$.

\textbf{19.} Equilíbrio $= 100 \times \frac{1.005^{60} - 1}
{1.005 - 1} \approx 100 \times \frac{0.3489}{0.005} \approx
6\,977$ euros --- dos quais $6\,000$ depositado e sobre $977$
merecido: \omterm{def:g11:seq:geometric}{geométrico} somas são a língua nativa do banco.

\textbf{20.} Fórmulas explícitas respondem ``o que é $ u_{1000}$''
instantaneamente; recorrências descrevem como os sistemas realmente evoluem ---
a arte é converter o segundo no primeiro. \omterm{def:g11:seq:arithmetic}{Aritmética}
\omterm{def:g11:seq:sequence}{sequências} adicionar, \omterm{def:g11:seq:geometric}{geométrico} multiplicam-se e cada família possui um
fórmula da soma (emparelhamento de Gauss; o truque da duplicação). O
truque de \omterm{pb:g10:functions:1}{ponto fixo} e auxiliar converte todos os afins
recorrência em um \omterm{def:g11:seq:geometric}{geométrico} um --- empréstimos, populações e
a torre caiu toda nela. Fibonacci não obedece a nenhuma família, ainda
identidades telescópicas capturavam suas somas e quadrados; está cheio
retrato (uma fórmula exata, o limite de ouro) aguarda mais forte
ferramentas.
\end{solution}
